\documentclass[11pt,a4paper]{article}
\usepackage[a4paper,margin=2.3cm]{geometry}\usepackage[T1]{fontenc}\usepackage[utf8]{inputenc}
\usepackage{parskip,hyperref,microtype,xcolor,needspace}\hypersetup{hidelinks}\setlength{\emergencystretch}{2em}
\begin{document}\begin{center}{\Large\bfseries Response to Reviewer 1 and Reviewer 3}\\[4pt]COMMT-D-26-00255 -- Computational Materials Today\end{center}
We thank the reviewers for their detailed comments. The original reviewer wording below is reproduced from the official review PDF, with line wrapping normalized. Quoted values describe the reviewed version; our responses retain the frozen, verified revision results. The revision withdraws unsupported LEA novelty/runtime, trained-GNN, and independent explanation-quality claims. No new major experiments were introduced in this final audit. Locations refer to the compiled clean manuscript: p. denotes the printed manuscript page, followed by the PDF page index. Supplementary pages are their PDF page indices.

\section*{Reviewer 1}
\Needspace{8\baselineskip}\subsection*{Comment 1}
\begin{quote}\small 1. The novelty of the proposed framework needs to be clarified. The current method mainly combines existing techniques, including retrieval-based recommendation, weighted ranking, diversity optimization, graph signals, and RAG. The authors should better explain the fundamental methodological contribution of QMOF-Rec compared with existing multi-objective recommendation and materials screening approaches.\end{quote}
\textbf{Response:} We agree that integration of established components alone should not be presented as a new optimization principle. The Introduction and four-item contribution list now identify the system-level formulation: query-conditioned screening under incomplete metadata, separate masked utility and material diversity, auditable portfolio trade-offs, and auxiliary explanation. We removed the LEA core-novelty and first-framework claims.

\textbf{Changes in revised manuscript:} Section 1, p. 1 (PDF page 2); Section 23, p. 29 (PDF page 30).

\Needspace{8\baselineskip}\subsection*{Comment 2}
\begin{quote}\small 2. The advantage of the Lotus Effect Algorithm (LEA) is not sufficiently demonstrated. The reported results show that LEA achieves similar ranking performance to several baselines, with only limited improvement in diversity. The authors should provide stronger comparisons with advanced optimization methods and discuss why LEA is preferable for this task.\end{quote}
\textbf{Response:} We added a SciPy differential-evolution portfolio baseline on identical pools, masks, and hybrid distances, with 30 members and 1,800 list-objective evaluations. It obtains Rel@5=0.8479 and diversity=0.0242. The fixed list objective and initialization are disclosed in Methods and the supplement; it is not asserted to reproduce LEA's history-dependent fitness. Revised Table 4 does not claim universal LEA preference. Moreover, top-five MMR takes 12.70 ms versus LEA 87.26 ms, with the same MMR selections. The previous speed-based preference is withdrawn.

\textbf{Changes in revised manuscript:} Section 14, p. 21 (PDF page 22); Table 4, p. 22 (PDF page 23); Section 20, p. 28 (PDF page 29); Supplementary Table S1, p. 1; Supplementary Table S3, p. 3.

\Needspace{8\baselineskip}\subsection*{Comment 3}
\begin{quote}\small 3. The evaluation mainly measures ranking consistency rather than actual materials discovery capability. The recommended candidates are evaluated using predefined objectives such as relevance, density, and band gap, but no external validation using experimental data, adsorption performance, or known high-performance MOFs is provided. Additional validation or elaboration is needed to demonstrate practical usefulness.\end{quote}
\textbf{Response:} We added scenario-level comparisons of observed density and band gap for the catalog, candidate pool, Random, WeightedSum, and LEA, with missingness and descriptor thresholds (supplementary descriptor table and raw screening.csv). These expose whether shortlist descriptors follow the stated screening intent. They use descriptors involved in scoring, so they are explicitly sanity checks, not independent materials-discovery validation. No adsorption uptake, porosity, synthesis, new DFT, or expert-label validation is claimed; that limitation remains unresolved by the available benchmark.

\textbf{Changes in revised manuscript:} Section 19, p. 27 (PDF page 28); Section 21, p. 29 (PDF page 30); Supplementary Table S5, p. 3.

\Needspace{8\baselineskip}\subsection*{Comment 4}
\begin{quote}\small 4. The graph-based representation should be described more carefully. The current graph features are derived from formula-level and metadata relationships rather than atomistic structures. The manuscript should avoid implying that these features capture genuine structural information and clearly distinguish them from CIF-based graph neural network approaches.\end{quote}
\textbf{Response:} The revised graph sections provide the actual formulas for proxies S and A and state that the benchmark neither constructs a graph nor loads trained GNN weights. The historical GraphSAGE-only/GAT-only labels are replaced by proxy scoring; these policies also retain 85\% relevance, so they were not graph-only. CIF files are present but unused in this protocol. The old trained-property prediction numbers and graph topology counts could not be reproduced and have been withdrawn; Table 6 now audits implementation scope.

\textbf{Changes in revised manuscript:} Section 8, p. 17 (PDF page 18); Section 9, p. 17 (PDF page 18); Table 6, p. 24 (PDF page 25); Table 2, p. 12 (PDF page 13); Figure 1, p. 9 (PDF page 10).

\Needspace{8\baselineskip}\subsection*{Comment 5}
\begin{quote}\small 5. The role of the LLM/RAG component is unclear. Currently, the LLM mainly provides explanation rather than contributing to candidate selection or optimization. The authors should clarify whether this component is a core scientific contribution or an auxiliary interface module.\end{quote}
\textbf{Response:} Explanation is now explicitly auxiliary and cannot alter candidate selection or ranking. Audit of the reproduced 30-template routine showed that several scores and all violation flags are directly assigned. We therefore replaced the nominal quality-score table with an implementation audit rather than describing perfect scores or zero flags as independent evidence. Live LLM, claim-checking, citation fidelity, and expert usefulness evaluation remain future work.

\textbf{Changes in revised manuscript:} Section 7, p. 16 (PDF page 17); Section 18, p. 27 (PDF page 28); Table 9, p. 27 (PDF page 28).

\Needspace{8\baselineskip}\subsection*{Comment 6}
\begin{quote}\small 6. More comprehensive baseline comparisons are needed. The evaluation should include modern AI-based materials screening methods, such as graph neural network-based ranking, transformer-based material representations, Bayesian optimization, or active learning approaches.\end{quote}
\textbf{Response:} DE portfolio addresses the request for a stronger optimization comparator; it does not replace a learned materials-representation comparator. Learned-representation comparisons remain an acknowledged limitation. We did not fabricate trained GNN or transformer comparison rows. A matched learned benchmark requires trained checkpoints, data splits, and comparable targets, which the final protocol does not supply; the available GAT module is not a trained or evaluated implementation within the final ranking protocol. Bayesian optimization and active learning require a sequential acquisition objective and oracle, absent here. The Discussion states this remaining scope limitation rather than claiming that the requested learned-model comparison is complete.

\textbf{Changes in revised manuscript:} Section 14, p. 21 (PDF page 22); Section 20, p. 28 (PDF page 29); Section 21, p. 29 (PDF page 30).

\Needspace{8\baselineskip}\subsection*{Comment 7}
\begin{quote}\small 7. The contribution of individual components requires further ablation analysis. The manuscript is recommended to provide clearer ablation studies to quantify the effects of LEA, diversity optimization, graph features, retrieval module, and missing-value handling.\end{quote}
\textbf{Response:} Revised Table 7 contains the LEA reference, gamma-only ablation, true no-diversity ablation, no-self-cleaning variant, proxy augmentation variants, a same-pool zero-imputation scoring diagnostic, and a uniform-candidate-pool diagnostic. Turning off both gamma and the redundancy penalty changes diversity from 0.0224 to 0.0187; the original gamma-only label was misleading. The uniform-pool row changes the feasible set and is not treated as a same-pool optimizer result. All coefficients and raw outputs are recorded.

\textbf{Changes in revised manuscript:} Section 16, p. 25 (PDF page 26); Table 7, p. 25 (PDF page 26); Section 17, p. 25 (PDF page 26); Supplementary Table S2, p. 2.

\Needspace{8\baselineskip}\subsection*{Comment 8}
\begin{quote}\small 8. The terminology should be refined. Since the framework mainly ranks existing QMOF database entries rather than discovering new materials, terms such as ``materials discovery'' should be used carefully. The work may be more accurately described as a materials screening or recommendation framework. If it is indeed truly generative, please highlight this.\end{quote}
\textbf{Response:} The evaluated task is now consistently framed as screening or prioritizing existing QMOF records. The Abstract and Introduction state that the system is not generative. Discovery terminology is retained only for the wider literature or future physical validation, and interface screenshots are not presented as evidence of novel materials.

\textbf{Changes in revised manuscript:} Section 1, p. 1 (PDF page 2); Section 14, p. 21 (PDF page 22); Table 5, p. 24 (PDF page 25); Section 23, p. 29 (PDF page 30).

\Needspace{8\baselineskip}\subsection*{Comment 9}
\begin{quote}\small 9. The computational efficiency should be reported in more detail. The authors should provide runtime and scalability analysis, especially for the LEA reranking process, and compare its computational cost with other ranking methods.\end{quote}
\textbf{Response:} The revision measures WeightedSum, TOPSIS, LEA, and top-five MMR for pools of 50, 100, 200, 500, and 1,000, with five scenarios, three seeds, and three repeats (45 measurements per method/size). Means, standard deviations, scope, and raw timings are provided. The original full-ranking MMR reproduced at 2209.67 ms, but its exactly matching top-five prefix needs 12.70 ms. We withdraw the approximately 25-fold LEA speed claim and discuss fixed-K complexity and excluded retrieval/generation costs.

\textbf{Changes in revised manuscript:} Figure 4, p. 23 (PDF page 24); Section 19, p. 27 (PDF page 28); Supplementary Table S3, p. 3; Supplementary Table S4, p. 3.

\Needspace{8\baselineskip}\subsection*{Comment 10}
\begin{quote}\small 10. The manuscript organization and contribution statements should be revised. The abstract is presented in multiple paragraphs, which is inconsistent with the common format of research articles in materials science journals. Moreover, the numbered contribution list contains several items describing general framework components or evaluation procedures rather than clear methodological advances. The authors should revise the presentation style and better distinguish genuine scientific contributions from implementation details.\end{quote}
\textbf{Response:} The abstract is now one paragraph. The contribution list is reduced to four scoped statements and separates the scientific formulation and empirical findings from general implementation capabilities. It incorporates the corrected ablation and runtime conclusions rather than repeating the previous stronger LEA claims.

\textbf{Changes in revised manuscript:} Abstract, PDF page 1; Section 1, p. 1 (PDF page 2).

\clearpage\section*{Reviewer 3}\subsection*{Original complete comment}
\begin{quote}\small Section 1 states that the "core novelty is the mask-aware post-retrieval reranking layer" (LEA). The reported results do not support LEA as an improvement over the baselines: (Table 4): LEA ties WeightedSum, TOPSIS and MMR on Rel@5 (0.8479) and NDCG@5 (1.0000). Note that Rel@K (Eq. 26) is the mean of the same score R (Eq. 6) that WeightedSum sorts by, so a tie with WeightedSum on these metrics is expected for any method that recovers its ordering; please state this and specify rel\_r used in Eq. 27. Now if we look at Table 4 that compares LEA with the other methods this is what we find: Diversity (Table 4): LEA (0.0224) is exceeded by ParetoCrowding (0.0263), MMR (0.0241), SemanticOnly (0.0240) and Random (0.0237). Ablation (Table 7): "LEA no diversity" and "LEA no self-clean" produce metrics identical to baseline LEA (0.8479 / 1.0000 / 0.0224), so LEA's distinguishing mechanisms have no measurable effect in this configuration. Per-scenario (Sec. 17): the mean LEA--WeightedSum diversity difference (0.0037) is driven almost entirely by the CO$_2$ scenario (0.0125); it is 0.0000 for photocatalysis. So instead of being the core novelty, a better framing would be, LEA is second on diversity and on the hypervolume proxy (0.0141 vs 0.0150 for MMR) while running \textasciitilde{}25$\times$ faster than MMR (83.92 vs 2114.36 ms). I recommend repositioning LEA as a relevance-preserving middle ground between WeightedSum and MMR rather than as the core novelty. To support the middle-ground framing, please add (i) per-scenario results for Table 4, and (ii) a sensitivity sweep over $\gamma$ (and the 0.05 redundancy-penalty weight) showing where, the diversity term becomes active. The SemanticOnly, WeightedSum, MMR and LEA rows appear verbatim in both Table 4 and Table 7. Please merge them into a single results table, or restrict Table 7 to the rows not already in Table 4 (LEA no diversity, LEA no self-clean, GraphSAGE only, GAT only, LEA + GraphSAGE, LEA + GAT), with a note that the LEA baseline row is repeated for reference.\end{quote}
\Needspace{7\baselineskip}\subsection*{Response issue 1: LEA contribution and novelty}
\textbf{Response:} We agree and removed that statement. The contribution is the auditable incomplete-metadata screening formulation. LEA is an evaluated policy with scenario-dependent behavior, not a newly established superior optimizer.

\textbf{Changes in revised manuscript:} Section 1, p. 1 (PDF page 2); Section 23, p. 29 (PDF page 30).

\Needspace{7\baselineskip}\subsection*{Response issue 2: Utility-aligned relevance and NDCG definition}
\textbf{Response:} The metrics section now defines rel\_r = R(m\_(r),q). IDCG sorts this same masked utility in the identical pool. WeightedSum maximizes both aligned metrics by construction; ties show utility preservation and are not independent evidence of relevance accuracy or physical discovery.

\textbf{Changes in revised manuscript:} Section 13.2, p. 20 (PDF page 21).

\Needspace{7\baselineskip}\subsection*{Response issue 3: Diversity ordering}
\textbf{Response:} The original values reproduce: ParetoCrowding 0.0263, MMR 0.0241, SemanticOnly 0.0240, Random 0.0237, LEA 0.0224, and WeightedSum/TOPSIS 0.0187. We state explicitly that LEA is not the most diverse. The added DE portfolio comparator reaches 0.0242, with its separate list objective disclosed.

\textbf{Changes in revised manuscript:} Table 4, p. 22 (PDF page 23); Supplementary Table S1, p. 1.

\Needspace{7\baselineskip}\subsection*{Response issue 4: Gamma-only versus true no-diversity ablation}
\textbf{Response:} The original row disabled gamma but retained the 0.05 exponential redundancy penalty. It is renamed gamma=0 only. Setting both coefficients to zero gives diversity 0.0187, versus baseline 0.0224. Gamma-only and no-self-cleaning aggregate null effects are preserved and discussed. We do not attribute the baseline improvement to both mechanisms independently.

\textbf{Changes in revised manuscript:} Section 16, p. 25 (PDF page 26); Table 7, p. 25 (PDF page 26).

\Needspace{7\baselineskip}\subsection*{Response issue 5: Complete per-scenario results}
\textbf{Response:} Supplementary Table S1 reports all seven original methods plus DE for all five scenarios. The reproduced LEA-minus-WeightedSum differences are 0.0125, 0.0000, 0.0028, 0.0029, and 0.0001. CO2 contributes about 68\% of their sum. Five scenarios are treated as descriptive blocks; ten seeds per scenario are not 50 independent scientific tasks.

\textbf{Changes in revised manuscript:} Section 17, p. 25 (PDF page 26); Supplementary Table S1, p. 1.

\Needspace{7\baselineskip}\subsection*{Response issue 6: Intermediate diversity and corrected runtime interpretation}
\textbf{Response:} The intermediate-diversity description is supported, but our additional matched-workload audit does not support the proposed runtime rationale. The original MMR computes all 100 ranks; stopping at five preserves every tested top-five list and reduces its mean runtime to 12.70 ms, below LEA's 87.26 ms. We report this correction candidly and do not retain a misleading LEA speed advantage.

\textbf{Changes in revised manuscript:} Table 4, p. 22 (PDF page 23); Figure 4, p. 23 (PDF page 24); Section 20, p. 28 (PDF page 29); Supplementary Table S3, p. 3; Supplementary Table S4, p. 3.

\Needspace{7\baselineskip}\subsection*{Response issue 7: Gamma sensitivity and observable activation}
\textbf{Response:} The revision sweeps gamma=0, 0.025, 0.05, 0.10, 0.20, 0.40, 0.80, jointly with five penalty coefficients. At penalty 0.05, gamma=0 retains the baseline sets in all 50 runs; gamma=0.80 changes 1/50 sets. We distinguish changes in sets, order, unrounded metrics, and recorded fitness-history hashes. The main sensitivity table and heatmap show where output effects appear, including null regions.

\textbf{Changes in revised manuscript:} Section 17, p. 25 (PDF page 26); Table 8, p. 26 (PDF page 27); Figure 5, p. 26 (PDF page 27); Supplementary Table S2, p. 2.

\Needspace{7\baselineskip}\subsection*{Response issue 8: Redundancy-coefficient sensitivity}
\textbf{Response:} The joint grid includes penalty coefficients 0, 0.025, 0.05, 0.10, and 0.20, for 35 settings and 1,750 query/seed runs. The supplement reports every pair, and sensitivity.csv includes scenario-level metrics and selected IDs. The baseline remains predefined; the grid is a robustness analysis and not outcome-driven retuning.

\textbf{Changes in revised manuscript:} Section 17, p. 25 (PDF page 26); Figure 5, p. 26 (PDF page 27); Supplementary Table S2, p. 2.

\Needspace{7\baselineskip}\subsection*{Response issue 9: Distinct content of Tables 4 and 7}
\textbf{Response:} Table 4 contains principal methods and DE. Table 7 contains component diagnostics, with only LEA repeated as the clearly identified reference row. SemanticOnly, WeightedSum, and MMR are no longer duplicated there; the zero-imputed WeightedSum row is a distinct labeled diagnostic.

\textbf{Changes in revised manuscript:} Table 4, p. 22 (PDF page 23); Table 7, p. 25 (PDF page 26).

\end{document}